\section{The regime \texorpdfstring{$\gamma<1$}{gamma < 1}: a frozen kernel and three posterior limits}
\label{sec:gamma-less-one}

Throughout this section the activation is $\tanh$, the depth $L\ge2$ and
input dimension $d\ge1$ are fixed, and the width tends to infinity. Write
$\mathcal X=\{-1,1\}^d$, $M=2^d$, and let $a>0$ be the required margin.
All limits below hold at each fixed $\gamma<1$; uniformity in dimension,
depth, or $\gamma\uparrow1$ is not asserted. Define
\begin{equation}\label{eq:arcsine-kernel}
 K^{(0)}(x,y)=\frac{1+x\cdot y}{d+1},\qquad
 K^{(s+1)}(x,y)=\frac2\pi\arcsin K^{(s)}(x,y),\qquad
 K=K^{(L-1)}.
\end{equation}
Thus the limiting kernel uses the saturated activation. It is not the
unit-variance, finite-gain $\tanh$ kernel.

\begin{theorem}[Exponential freezing of the feature kernel]
\label{thm:kernel-freezing}
Let $h_i^s(x)$ denote the activation of hidden neuron $i$ in layer $s$,
and set
\[
 Q_N^s(x,y)=\frac1N\sum_{i=1}^N h_i^s(x)h_i^s(y).
\]
For every $\epsilon>0$ and each $1\le s\le L-1$, there are constants
$c,C>0$ and $N_0$ such that
\[
 \Pr\!\left\{\max_{x,y}|Q_N^s(x,y)-K^{(s)}(x,y)|>\epsilon\right\}
 \le C e^{-cN}\qquad(N\ge N_0).
\]
The matrix $K$ is positive definite. Conditional on all hidden layers,
the output vector $Z_N=(z_N(x))_{x\in\mathcal X}$ has the exact law
\begin{equation}\label{eq:conditional-output-kernel}
 Z_N\mid (\text{hidden layers})\sim
 \mathcal N(0,N^{-\gamma}Q_N^{L-1}).
\end{equation}
Consequently $N^{\gamma/2}Z_N$ converges in distribution to
$G\sim\mathcal N(0,K)$.
\end{theorem}

\begin{proof}
Put $b_N=N^{(1-\gamma)/2}\to\infty$. Conditional on the preceding
layer, different preactivation rows are independent Gaussian vectors
with covariance $b_N^2Q_N^{s-1}$. For the first layer, replace this
covariance matrix by the deterministic matrix
$b_N^2R_0$, where $R_0(x,y)=(1+x\cdot y)/d$; its correlation matrix is
$K^{(0)}$.

If $U\sim\mathcal N(0,q)$, then bounding the Gaussian density by its
maximum and integrating $1-\tanh v$ gives
\[
 \mathbb E|\tanh(bU)-\operatorname{sign}U|
 \le \frac{2\log2}{b\sqrt{2\pi q}}.
\]
For a centered Gaussian pair with correlation $r$, the expected product
of its signs is $2\arcsin(r)/\pi$. Thus, uniformly when coordinate
variances are bounded below, the expected product of the two saturated
activations differs from this arcsine expression by $O(b_N^{-1})$.
The appended constant input implies
$|K^{(s)}(x,y)|<1$ for distinct $x,y$. In a sufficiently small fixed
neighborhood of $K^{(s)}$, the map sending a covariance matrix to its
normalized arcsine matrix is Lipschitz, and all diagonal variances are
bounded below. Finally, conditional Hoeffding bounds for the $N$
independent products in $[-1,1]$ give
\[
 \Pr\!\left\{\max_{x,y}|Q_N^s(x,y)
 -\mathbb E[Q_N^s(x,y)\mid\text{preceding layer}]|>t
 \,\middle|\,\text{preceding layer}\right\}
 \le 2M^2 e^{-Nt^2/2}.
\]
Induction over the fixed number of layers proves the concentration
statement, absorbing the deterministic $O(b_N^{-1})$ error by taking
$N$ sufficiently large.

For positive definiteness, $K^{(1)}$ is the Gram matrix of the functions
$u\mapsto\operatorname{sign}(u\cdot(1,x))$ under a nondegenerate
Gaussian measure. They are linearly independent: in a vanishing linear
combination, cross the hyperplane $u\cdot(1,x)=0$ at a point on none of
the other finitely many hyperplanes. Only the corresponding term
changes sign, so its coefficient is zero. Such a point exists because
the augmented vectors $(1,x)$ are pairwise nonparallel. The positive
power-series coefficients of the arcsine function and the Schur product
theorem then imply $K^{(s+1)}\succeq(2/\pi)K^{(s)}$, preserving strict
positive definiteness. The independent Gaussian readout proves
\eqref{eq:conditional-output-kernel}, and convergence follows.
\end{proof}

\subsection{Positive \texorpdfstring{$\gamma$}{gamma}: the actual posterior is a kernel-energy learner}

For a Boolean function $g$ define
\begin{equation}\label{eq:kernel-orthant-energy}
 E_K(g)=\min\left\{\frac12z^TK^{-1}z:
               g(x)z(x)\ge a\text{ for every }x\in\mathcal X\right\}.
\end{equation}
This minimum is finite and has a unique minimizing vector, by strict
convexity and coercivity.

\begin{theorem}[Kernel posterior limit]
\label{thm:kernel-posterior}
For each fixed $0<\gamma<1$, $Z_N$ satisfies a large-deviation principle
with speed $N^\gamma$ and good rate function
$I(z)=z^TK^{-1}z/2$. In particular,
\begin{equation}\label{eq:kernel-orthant-rate}
 \lim_{N\to\infty}-N^{-\gamma}
 \log\Pr\{g(x)Z_N(x)>a\text{ for every }x\}=E_K(g).
\end{equation}
Let $\mathcal G_S$ be the Boolean functions consistent with the observed
training labels and let $E_S=\min_{h\in\mathcal G_S}E_K(h)$. Under the
all-input-margin prior followed by noiseless sign conditioning,
\begin{equation}\label{eq:kernel-posterior-rate}
 \lim_{N\to\infty}-N^{-\gamma}
 \log\Pr\{\operatorname{sign}Z_N=g\mid\text{all margins and }S\}
 =E_K(g)-E_S\qquad(g\in\mathcal G_S).
\end{equation}
The posterior therefore concentrates on the functions in
$\operatorname*{argmin}_{g\in\mathcal G_S}E_K(g)$.
\end{theorem}

\begin{proof}
On a common probability space let
\[
 Z_N=N^{-\gamma/2}(Q_N^{L-1})^{1/2}\xi,
 \qquad Y_N=N^{-\gamma/2}K^{1/2}\xi,
 \qquad \xi\sim\mathcal N(0,I_M),
\]
with $\xi$ independent of the hidden layers. The square-root map is
locally Lipschitz near the positive definite $K$. For fixed
$\delta,\epsilon>0$, Theorem~\ref{thm:kernel-freezing} consequently gives
\[
 \Pr\{\|Z_N-Y_N\|>\delta\}
 \le C_\epsilon e^{-c_\epsilon N}
   +\Pr\!\left\{\|\xi\|>
          \frac{\delta N^{\gamma/2}}{C\epsilon}\right\}.
\]
Divide logarithms by $N^\gamma$, first let $N\to\infty$, and then let
$\epsilon\downarrow0$. Because $\gamma<1$, the resulting upper limit
is $-\infty$. Thus $Z_N$ is exponentially equivalent, at speed
$N^\gamma$, to the elementary small-noise Gaussian family $Y_N$.
The latter has rate $I$, proving the large-deviation statement.

The strict margin orthant is open and its closure is the feasible set
in~\eqref{eq:kernel-orthant-energy}. Scaling its unique minimizer by
$1+\eta$ puts it in the open orthant and approaches the same energy
as $\eta\downarrow0$. The open lower bound and closed upper bound
therefore match. Summing over the finitely many compatible sign
patterns and subtracting logarithmic asymptotics proves
\eqref{eq:kernel-posterior-rate}.
\end{proof}

If the minimizing Boolean function is unique, the limiting posterior
is its point mass. If there is a tie, the exponential rates do not
determine the mixture weights. The result does not prescribe a uniform
mixture of minimizers.

The limit is governed by a fixed kernel, but its optimization problem
retains the constraint $|z(x)|\ge a$ at every unseen input. Its feasible
set is
\[
 \{z:f(x_i)z(x_i)\ge a\text{ at observed inputs},\quad
             |z(x)|\ge a\text{ at every input}\}.
\]
This is a nonconvex union of orthants. Dropping the unseen-input
constraints produces a different learner. In particular, ordinary
kernel interpolation, Gaussian-process regression, and the usual
hard-margin support-vector machine do not describe this posterior.

\subsection{Zero and negative \texorpdfstring{$\gamma$}{gamma}}

\begin{proposition}[The remaining fixed-kernel limits]
\label{prop:nonpositive-gamma}
At $\gamma=0$, the posterior on Boolean functions converges to that of
$G\sim\mathcal N(0,K)$ conditioned on $|G(x)|>a$ at every input and on
the observed training signs. At $\gamma<0$, it converges to the sign
law of the same Gaussian vector conditioned only on observed training
signs.
\end{proposition}

\begin{proof}
The normalized outputs converge weakly to the nondegenerate $G$.
At $\gamma=0$, every relevant conditioning event has boundary in a
finite union of hyperplanes of Gaussian probability zero, and its
probability is strictly positive. Conditional probabilities converge.
At $\gamma<0$, the normalized margin is
$aN^{\gamma/2}\to0$. The probability that any normalized coordinate
falls within this margin tends to zero. Every prescribed finite set
of observed signs has positive limiting probability, so the same
conditional convergence argument applies.
\end{proof}

\subsection{What the kernel limit says about parity}

Cube symmetries diagonalize $K$ in the parity basis. If
$\chi_U(x)=\prod_{i\in U}x_i$ and $\lambda_U$ is its eigenvalue for
the normalized operator $M^{-1}K$, then
\begin{equation}\label{eq:parity-kernel-energy}
 E_K(\chi_U)=\frac{a^2}{2\lambda_U}.
\end{equation}
Indeed $a\chi_U$ is feasible, and its gradient
$K^{-1}(a\chi_U)$ is a positive multiple of $\chi_U$. Its inner product
with every feasible displacement is nonnegative; the first-order
criterion for convex minimization proves optimality.

Equation~\eqref{eq:parity-kernel-energy} is an exact prior-rate
calculation. It is not a sample-complexity lower bound. For
$0<\gamma<1$, that question requires comparison with the energies of
all compatible Boolean functions; for $\gamma=0$ it concerns globally
margin-conditioned Gaussian orthant probabilities; for $\gamma<0$
it concerns zero-margin Gaussian sign probabilities. A small parity
eigenvalue by itself does not answer any of these posterior-learning
questions. The regression subspace argument also does not apply: the
full-cube conditioning need not leave the posterior mean in the span
of training kernel sections. Moreover a lower bound stated in the
parity size $k$ must specify the ambient dimension, depth, input
distribution, order of limits, and whether error means posterior
misclassification probability or majority-vote error.
